% META: {"id": "TEXP-GRA-CO-007", "chapitre": "TEXP-GRAPHES", "type_objet": "corrige", "exercice_ref": "TEXP-GRA-EX-007", "capacites_codes": ["C2"], "sources_inspiration": [], "status": "generated"}
% BEGIN-VERIFY
% import collections
% aretes = [("U", "R"), ("U", "L"), ("R", "C"), ("L", "C"), ("C", "P")]
% voisins = collections.defaultdict(set)
% for u, v in aretes:
%     voisins[u].add(v); voisins[v].add(u)
% degre = {s: len(v) for s, v in voisins.items()}
% assert degre == {"U": 2, "R": 2, "L": 2, "C": 3, "P": 1}
% assert max(degre, key=degre.get) == "C"
% distance = {"U": 0}
% file = collections.deque(["U"])
% while file:
%     s = file.popleft()
%     for t in voisins[s]:
%         if t not in distance:
%             distance[t] = distance[s] + 1
%             file.append(t)
% assert distance["P"] == 3
% END-VERIFY
\begin{corrige}{TEXP-GRA-EX-007}
\textbf{1.} Les sommets sont les cinq services ; une arête relie deux services
entre lesquels une liaison directe existe. Le graphe est non orienté, car une
liaison fonctionne dans les deux sens.

\textbf{2.} Oui : Urgences--Radiologie--Chirurgie--Pharmacie, soit trois
étapes. Aucune chaîne plus courte n'existe, la Pharmacie n'étant reliée qu'à
la Chirurgie, elle-même non adjacente aux Urgences.

\textbf{3.} Les degrés valent $2$ pour Urgences, Radiologie et Laboratoire,
$1$ pour Pharmacie, et $3$ pour Chirurgie. La Chirurgie est le sommet de degré
maximal : c'est le service le plus sollicité comme intermédiaire.
\end{corrige}
